\section{Dataset generation and the runtime qualification boundary}
\label{sec:b-data}

\subsection{Collection sequence}

One root is generated as follows:
\begin{enumerate}[label=(\arabic*),leftmargin=2em]
\item Run the prefix once to the declared checkpoint under the frozen agent.
\item Quiesce execution: finish or cancel in-flight tools according to the
  contract.
\item Capture the live declared state and write a manifest (state surface,
  hashes, versions, budget counters).
\item Create isolated sibling states immediately.
\item Randomize or counterbalance branch execution order by a predeclared rule.
\item Execute $g_0$ and each feasible intervention for $R$ repetitions each.
\item Resume the same $\rho$ under the registered horizon and budget semantics.
\item Adjudicate final artifacts with a hidden, regime-blind evaluator.
\item Write append-only records and branch accounting.
\end{enumerate}

Live forking reduces historical reconstruction requirements; it does not prove
equal effective state.  If the runtime only supports restarting from a
serialized transcript, the operation is described as \emph{re-instantiation},
and the interface is validated as such; it is not called an exact process fork.

\subsection{Declared state surface}

The outcome-relevant subset of the following is checked and versioned:
workspace files, permissions and artifact versions; tool/process state,
subprocesses and job IDs; environment and library versions; history, prompt,
model identity and decoding settings; RNG state and the seed-coupling rule;
budget counters, elapsed time and remaining horizon; caches, mutable shared
stores, symlink targets and provider sessions; and network/external services if
allowed.  A pure output-neutral cache need not invalidate causal outcomes merely
because it is warm; its effect is tested on both the outcomes and the chosen
cost endpoint, and if wall-time budget makes cache warmth outcome-relevant it is
included in the contract.  Reference validity cannot be established by
comparing two implementations that share the same hidden defect; the
independently justified reference boundary is stated with the contract.

\subsection{The qualification gate (B-Q)}

A separate pool with engineering labels is used.  Required checks: identical
serialized starting payloads where equality is claimed; immutable read-only
inputs and isolated writable state; no cross-branch mutable links, queues or
counters; pre-treatment mask and feature construction; stable
action/continuation/adjudicator versions; branch-order and seed-handling checks;
deterministic known-answer controls and stochastic-law diagnostics where
relevant; and explicit missing-state and tool-failure handling.  Tolerances are
matched to the state contract before the scientific pilot; arbitrary determinism
is not required from a genuinely stochastic model.  An engineering qualification
check is evidence about declared components, not proof of all hidden
dependencies: unresolved outcome-relevant state blocks causal wording for that
substrate.  Only after B-Q passes for a specific runtime/configuration, with a
bounded statement of what was checked, may new pilot branch contrasts be
interpreted scientifically.

\subsection{Branch record schema}

Every branch record carries: schema version, paper, study, protocol hash, run
id, template id, root id, checkpoint id and rule version, split, evaluator-only
condition label, runtime manifest hash, snapshot and history hashes, feature
registry version, pre-treatment features and mask, availability evidence hash,
action id and version, continuation version, repetition id, seed-coupling id,
branch order, the outcome fields (success, utility, invalid), the cost vector,
execution status and failure reason, adjudicator version, and artifact hashes.
The learner dataset is a restricted projection of this record.  The outcome
table and the evaluator-only metadata are not reachable through feature
construction, and the projection is validated by a leak test that must fail
when the projection is given access to either.

\subsection{Splits}

Qualification and pilot pools are disjoint from final evaluation.  Training fits
representations and learners; validation selects architecture, hyperparameters
and candidate policies; calibration fixes conservative thresholds and selection;
the final test evaluates frozen policies.  Splitting is at the template level:
multiple roots, seeds, fault variants or tasks generated from one template stay
in one split, and shared task datasets are recorded at that level too.
